
\[
  \Projz\bigl(F_k(\underline{L}^{\bullet})\bigr)\longrightarrow \underline{L}^{\bullet},\qquad ,
\]
et c'est un
isomorphisme sur chaque composante puisque $\underline{L}^{\bullet}$ a été supposé strictement parfait, d'où le ii).

2.7. Sur $\mathcal C$, le foncteur qui à un $\underline{B}$-Module gradué associe sa $n$-ième composante est un foncteur exact à valeurs dans $\mathcal C'$; il définit donc un foncteur :
\begin{equation}\tag{2.7.1}
  \operatorname{Comp}_n:D^+(\mathcal C)\longrightarrow D^+(\mathcal C').
\end{equation}
On peut, pour tout $n$, considérer le foncteur $\mathcal C'\to\mathcal C$ défini par la tensorisation par $\underline{B}(-n)$. L'homomorphisme (1.2.1) définit un morphisme de foncteurs sur $\mathcal C$, d'où un morphisme pour les foncteurs correspondants entre catégories dérivées :
\begin{equation}\tag{2.7.2}
  \operatorname{Comp}_n\otimes_{\underline{O}_S} \underline{B}(-n)
  \longrightarrow \operatorname{Id}_{D^+(\mathcal C)}.
\end{equation}
Soit $\underline{M}^{\bullet}$ un complexe $\mathcal C_0$-parfait. Nous allons montrer qu'il existe une suite de triangles distingués :
\[
\begin{array}{ccccccccccc}
& \underline{P}_1^{\bullet} &&&& \underline{P}_2^{\bullet} &&&& \underline{P}_k^{\bullet} & \\
\swarrow && \nwarrow && \swarrow && \nwarrow && \swarrow && \nwarrow \\
\underline{N}_1^{\bullet} & \longrightarrow & \underline{M}^{\bullet}, & \qquad &
\underline{N}_2^{\bullet} & \longrightarrow & \underline{P}_1^{\bullet}, & ...., &
\underline{N}_k^{\bullet} & \longrightarrow & \underline{P}_{k-1}^{\bullet},
\end{array}
\]
où $\underline{P}_k^{\bullet}$ est acyclique, et les $\underline{N}_i^{\bullet}$ du type \(\underline{K}_i^{\bullet}\otimes_{\underline{O}_S}\underline{B}(d)\), les $\underline{K}_i^{\bullet}$ étant des complexes parfaits au sens ordinaire sur $S$.

Puisque $\underline{M}^{\bullet}$ est $\mathcal C_0$-parfait, et que $S$ est quasi-compact, il existe un recouvrement fini de $S$, et un entier $d_0$ tel que sur tout ouvert de ce recouvrement, $\underline{M}^{\bullet}$ soit isomorphe à un complexe borné à termes localement libres de type fini sur $\underline{B}$, soit $\underline{M}_U'{}^{\bullet}$, tel que $\operatorname{Comp}_d(\underline{M}_U'{}^{\bullet})=0$ pour $d<d_0$.

